\documentclass[12pt]{article}

\usepackage[T1]{fontenc}
\usepackage[utf8]{inputenc}
\usepackage[brazil]{babel}
\usepackage{lmodern}
\usepackage{microtype}
\usepackage[a4paper,margin=2.5cm]{geometry}
\usepackage{setspace}
\usepackage{indentfirst}

\setstretch{1.15}

\title{Depois da Macroeconomia Keynesiana \thanks{Texto Traduzido livremente por Luiz Mario Andrade com o auxílio da ferramenta Chat GPT 5.4}}
\author{Robert E. Lucas e Thomas J. Sargent}
\date{1978}

\begin{document}

\maketitle

\section{Introdução}

Para o economista aplicado, a aplicação confiante e aparentemente bem-sucedida dos princípios keynesianos à política econômica ocorrida nos Estados Unidos na década de 1960 foi um evento de incomparável significado e satisfação. Esses princípios levaram a um conjunto de relações quantitativas simples entre a política fiscal e a atividade econômica em geral, cuja lógica básica podia ser (e foi) explicada ao público em geral, e que podia ser aplicada para produzir melhorias no desempenho econômico que beneficiavam a todos. Parecia uma economia tão livre de dificuldades ideológicas quanto, digamos, a química ou a física aplicada, prometendo uma expansão direta das possibilidades econômicas. Poder-se-ia discutir sobre como esse ganho inesperado deveria ser distribuído, mas parecia um simples lapso de lógica opor-se ao próprio ganho. Compreensivelmente e corretamente, essa promessa foi recebida inicialmente com ceticismo por não-economistas; a prosperidade de crescimento suave dos anos Kennedy-Johnson contribuiu muito para diminuir essas dúvidas.

Demos ênfase a esses dias idílicos da economia keynesiana porque, sem um esforço consciente, eles são difíceis de recordar hoje. Na presente década, a economia dos EUA passou por sua primeira grande depressão desde a década de 1930, acompanhada por taxas de inflação superiores a 10 por cento ao ano. Esses eventos foram transmitidos (com o consentimento dos governos envolvidos) a outros países avançados e, em muitos casos, foram amplificados. Esses eventos não surgiram de uma reversão reacionária a princípios "clássicos" ultrapassados de moeda apertada e orçamentos equilibrados. Pelo contrário, foram acompanhados por déficits orçamentários governamentais massivos e altas taxas de expansão monetária: políticas que, embora tivessem um risco admitido de inflação, prometiam, de acordo com a doutrina keynesiana moderna, um rápido crescimento real e baixas taxas de desemprego.

Que essas previsões estavam totalmente incorretas e que a doutrina na qual se baseavam é fundamentalmente falha, são agora simples questões de fato, não envolvendo novidades na teoria econômica. A tarefa que se impõe aos estudantes contemporâneos do ciclo econômico é a de fazer uma triagem dos destroços, determinando quais características desse notável evento intelectual chamado Revolução Keynesiana podem ser salvas e bem aproveitadas, e quais outras devem ser descartadas. Embora esteja longe de ser claro qual será o resultado desse processo, já é evidente que envolverá necessariamente a reabertura de questões básicas na economia monetária que têm sido vistas desde os anos trinta como "encerradas", e a reavaliação de cada aspecto da estrutura institucional dentro da qual a política monetária e fiscal é formulada nos países avançados.

Este artigo tem o caráter de um relatório de progresso inicial sobre esse processo de reavaliação e reconstrução. Começamos revisando a estrutura econométrica por meio da qual a teoria keynesiana evoluiu de "conversas" qualitativas desconexas sobre a atividade econômica para um sistema de equações que podia ser comparado aos dados de forma sistemática e fornecer um guia operacional na tarefa necessariamente quantitativa de formular políticas monetárias e fiscais. Em seguida, identificamos os aspectos dessa estrutura que foram centrais para o seu fracasso na década de setenta. Ao fazê-lo, nossa intenção será estabelecer que as dificuldades são fatais: que os modelos macroeconômicos modernos não têm valor para orientar a política econômica e que essa condição não será remediada por modificações em qualquer linha que esteja sendo seguida atualmente.

Este diagnóstico, se for bem-sucedido, sugerirá certos princípios que uma teoria útil dos ciclos econômicos deve possuir. Na parte final deste artigo, revisaremos algumas pesquisas recentes que são consistentes com esses princípios.

\section{Modelos Macroeconométricos}

A Revolução Keynesiana foi, na forma em que triunfou nos Estados Unidos, uma revolução de método. Isso não era a intenção de Keynes [13], nem é a visão de todos os seus seguidores mais eminentes. No entanto, se não virmos a revolução dessa maneira, é impossível dar conta de algumas de suas características mais importantes: a evolução da macroeconomia em uma disciplina quantitativa e científica, o desenvolvimento de descrições estatísticas explícitas do comportamento econômico, a crescente dependência de funcionários do governo em relação à perícia técnica econômica e a introdução do uso da teoria do controle matemático para gerenciar uma economia. É o fato de que a teoria keynesiana se prestou tão prontamente à formulação de modelos econométricos explícitos que explica a posição científica dominante que alcançou na década de 1960.

Como consequência disso, não há esperança de entender o sucesso da Revolução Keynesiana ou o seu eventual fracasso apenas no nível verbal em que o próprio Keynes escreveu. Será necessário conhecer algo sobre a maneira como os modelos macroeconométricos são construídos e as características que devem ter para "funcionar" como auxílios na previsão e avaliação de políticas. Para discutir essas questões, introduzimos alguma notação.

Um modelo econométrico é um sistema de equações envolvendo um número de variáveis endógenas (variáveis determinadas pelo modelo), variáveis exógenas (variáveis que afetam o sistema, mas não são afetadas por ele) e choques estocásticos ou aleatórios. A ideia é usar dados históricos para estimar o modelo e, em seguida, utilizar a versão estimada para obter estimativas das consequências de políticas alternativas. Por razões práticas, é comum usar um modelo linear padrão, assumindo a forma estrutural\footnote{A linearidade é uma questão de conveniência, não de princípio. Veja a Seção 6.3, abaixo.}:

\begin{equation}
A_0 y_t + A_1 y_{t-1} + \dots + A_m y_{t-m} = B_0 x_t + B_1 x_{t-1} + \dots + B_n x_{t-n} + \epsilon_t
\end{equation}
\begin{equation}
R_0 \epsilon_t + R_1 \epsilon_{t-1} + \dots + R_r \epsilon_{t-r} = u_t, \quad R_0 \equiv I.
\end{equation}

Aqui, $y_t$ é um vetor $(L \times 1)$ de variáveis endógenas, $x_t$ é um vetor $(K \times 1)$ de variáveis exógenas, e $\epsilon_t$ e $u_t$ são vetores $(L \times 1)$ de distúrbios aleatórios. As matrizes $A_j$ são cada uma $(L \times L)$; as $B_j$ são $(L \times K)$, e as $R_j$ são cada uma $(L \times L)$. O processo de distúrbio $(L \times 1)$ $u_t$ é assumido como um processo serialmente não correlacionado com $E u_t = 0$ e com matriz de covariância contemporânea $E u_t u_t' = \Sigma$ e $E u_t u_s' = 0$ para todo $t \neq s$. A característica definidora das variáveis exógenas $x_t$ é que elas não são correlacionadas com os $\epsilon$'s em todos os defasamentos, de modo que $E u_t x_s'$ é uma matriz $(L \times K)$ de zeros para todos os $t$ e $s$.

As equações (1) são $L$ equações nos $L$ valores atuais $y_t$ das variáveis endógenas. Cada uma dessas equações estruturais é uma relação comportamental, identidade ou condição de equilíbrio de mercado, e cada uma, em princípio, pode envolver uma série de variáveis endógenas. As equações estruturais geralmente não são "equações de regressão"\footnote{Uma "equação de regressão" é uma equação para a qual a aplicação de mínimos quadrados ordinários produzirá estimativas consistentes.} porque os $\epsilon_t$ supõe-se serem em geral, pela lógica do modelo, correlacionados com mais de um componente do vetor $y_t$ e muito possivelmente com um ou mais componentes dos vetores $y_{t-1}, \dots, y_{t-m}$.

O modelo estrutural (1) e (2) pode ser resolvido para $y_t$ em termos de $y$'s e $x$'s passados e choques passados. Este sistema de "forma reduzida" é:

\begin{equation}
y_t = - P_1 y_{t-1} - \dots - P_{r+m} y_{t-r-m} + Q_0 x_t + \dots + Q_{r+n} x_{t-n-r} + A_0^{-1} u_t
\end{equation}
onde\footnote{Nestas expressões para $P_s$ e $Q_s$, considere matrizes não definidas anteriormente (por exemplo, qualquer uma com subscritos negativos) como sendo zero.}
\[ P_s = A_0^{-1} \sum_{j=-\infty}^{\infty} R_j A_{s-j} \]
\[ Q_s = A_0^{-1} \sum_{j=-\infty}^{\infty} R_j B_{s-j} \]

As equações de forma reduzida são "equações de regressão", ou seja, o vetor de distúrbio $A_0^{-1} u_t$ é ortogonal a $y_{t-1}, \dots, y_{t-r-m}, x_t, \dots, x_{t-n-r}$. Isso decorre das suposições de que os $x$'s são exógenos e que os $u$'s são serialmente não correlacionados. Portanto, sob condições gerais, a forma reduzida pode ser estimada consistentemente pelo método dos mínimos quadrados. Os parâmetros populacionais da forma reduzida (3), juntamente com os parâmetros de uma autorregressão vetorial para $x_t$,

\begin{equation}
x_t = C_1 x_{t-1} + \dots + C_p x_{t-p} + a_t
\end{equation}
onde $E a_t = 0$ e $E a_t \cdot x_{t-j}' = 0$ para $j \geq 1$, descrevem completamente todos os primeiros e segundos momentos do processo $(y_t, x_t)$. Dado séries temporais suficientemente longas, boas estimativas dos parâmetros da forma reduzida — os $P_j$ e $Q_j$ — podem ser obtidas pelo método dos mínimos quadrados. Estimativas confiáveis desses parâmetros são tudo o que o exame dos dados por si só pode fornecer.

Em geral, não é possível trabalhar de trás para frente a partir das estimativas dos $P$ e $Q$ isoladamente para derivar estimativas únicas dos parâmetros estruturais, os $A_j, B_j$ e $R_j$. Em geral, um número infinito de matrizes $A, B$ e $R$ é compatível com um único conjunto de $P$ e $Q$. Este é o "problema de identificação" da econometria. Para derivar um conjunto de parâmetros estruturais estimados, é necessário saber muito sobre eles com antecedência. Se for imposta informação prévia suficiente, é possível extrair estimativas dos $(A_j, B_j, R_j)$ implícitas nos dados em combinação com a informação prévia.

Para fins de previsão \textit{ex ante}, ou a predição incondicional do vetor $y_{t+1}, y_{t+2}, \dots$ dadas observações de $y_s$ e $x_s, s \leq t$, a forma reduzida estimada (3), juntamente com (4), é suficiente. Isso é simplesmente um exercício de um tipo sofisticado de extrapolação, não exigindo compreensão dos parâmetros estruturais ou, o que quer dizer, da economia do modelo.

Para fins de previsão \textit{condicional}, ou a predição do comportamento futuro de alguns componentes de $y_t$ e $x_t$ condicionados a valores particulares de outros componentes selecionados pela política econômica, é necessário conhecer os parâmetros estruturais. Isso ocorre porque uma mudança na política necessariamente altera alguns dos parâmetros estruturais (por exemplo, aqueles que descrevem o comportamento passado das próprias variáveis de política) e, portanto, afeta os parâmetros da forma reduzida de maneira altamente complexa (veja as equações que definem $P_s$ e $Q_s$ abaixo de (3)). Sem o conhecimento de quais parâmetros estruturais permanecem invariantes à medida que a política muda e quais mudam (e como), um modelo econométrico não tem valor para avaliar políticas alternativas. Deve ficar claro que isso é verdade independentemente de quão bem (3) e (4) se ajustem aos dados históricos ou de quão bem funcionem na previsão incondicional.

Nossa discussão até este ponto tem sido em um alto nível de generalidade, e as considerações formais que revisamos não são de forma alguma específicas aos modelos keynesianos. O problema de identificar um modelo estrutural a partir de uma coleção de séries temporais econômicas é um problema que deve ser resolvido por qualquer pessoa que afirme ser capaz de dar conselhos econômicos quantitativos. Os modelos keynesianos mais simples são tentativas de solução para este problema, assim como as versões em larga escala atualmente em uso. O mesmo ocorre com os modelos monetaristas que implicam a desejabilidade de regras fixas de crescimento monetário. O mesmo se aplica, inclusive, aos conselhos de gabinete dados por economistas que afirmam estar fora da tradição econométrica, embora, neste caso, a estrutura implícita subjacente não seja exposta à crítica profissional. Qualquer procedimento que leve do estudo do comportamento econômico observado para a avaliação quantitativa de políticas econômicas alternativas envolve as etapas, executadas bem ou mal, explícita ou implicitamente, que delineamos acima.

\section{Macroeconometria Keynesiana}

Nos modelos macroeconométricos keynesianos, os parâmetros estruturais são identificados pela imposição de vários tipos de restrições \textit{a priori} nos $A_j, B_j$ e $R_j$. Essas restrições geralmente se enquadram em uma das seguintes categorias\footnote{Essas três categorias certamente não esgotam o conjunto de possíveis restrições de identificação, mas nos modelos macroeconométricos keynesianos a maioria das restrições de identificação recai em uma destas três categorias. Outros tipos possíveis de restrições de identificação incluem, por exemplo, conhecimento \textit{a priori} sobre componentes de $\Sigma$ e restrições entre equações entre elementos de $A_j, B_j$ e $C_j$. Nenhum desses últimos tipos de restrições é amplamente utilizado na macroeconometria keynesiana.}:

\begin{enumerate}
    \item[(a)] Fixação \textit{a priori} de muitos dos elementos das matrizes $A_j$ e $B_j$ como zero.
    \item[(b)] Restrições nas ordens de correlação serial e na extensão da correlação serial cruzada do vetor de distúrbios $\epsilon_t$, restrições que equivalem a fixar \textit{a priori} muitos elementos dos $R_j$ como zero.
    \item[(c)] Categorização \textit{a priori} de variáveis em "exógenas" e "endógenas". Uma abundância relativa de variáveis exógenas auxilia na identificação.
\end{enumerate}

Os grandes modelos macroeconométricos keynesianos existentes estão abertos a sérios desafios pela maneira como introduziram cada categoria de restrição.

A \textit{Teoria Geral} de Keynes foi rica em sugestões para restrições do tipo (a). Propôs uma teoria da determinação da renda nacional construída a partir de várias relações simples, cada uma envolvendo apenas algumas variáveis. Uma delas, por exemplo, era a "lei fundamental" que relacionava as despesas de consumo à renda. Isso sugeria uma "linha" nas equações (1) envolvendo o consumo atual, a renda atual e nenhuma outra variável, impondo assim muitas restrições de zero nos $A_i$ e $B_j$. Da mesma forma, a relação de preferência pela liquidez expressava a demanda por moeda como função apenas da renda e de uma taxa de juros. Ao traduzir os blocos de construção do sistema teórico keynesiano em equações explícitas, modelos da forma (1) e (2) foram construídos com muitas restrições teóricas do tipo (a).

Restrições nos coeficientes $R_i$ que governam o comportamento dos "termos de erro" em (1) são mais difíceis de motivar teoricamente, sendo os "erros", por definição, movimentos nas variáveis que a teoria econômica não pode explicar. Os primeiros econometristas adotaram premissas "padrão" de livros didáticos de estatística, restrições que se mostraram úteis na experimentação agrícola, a qual forneceu o principal impulso para o desenvolvimento da estatística moderna. Novamente, essas restrições, bem motivadas ou não, envolvem a definição de muitos elementos nos $R_j$ como iguais a zero, auxiliando na identificação da estrutura do modelo.

A classificação de variáveis em "exógenas" e "endógenas" também foi feita com base em considerações prévias. Em geral, as variáveis eram classificadas como "endógenas" quando eram, por um fato institucional, determinadas em grande parte pelas ações de agentes privados (como o consumo ou despesas de investimento privado). Variáveis exógenas eram aquelas sob controle governamental (como taxas de impostos ou a oferta monetária). Essa divisão pretendia refletir o significado comum da palavra "endógena" como "determinada pelo sistema [econômico]" e "exógena" como "afetando o sistema [econômico], mas não sendo afetada por ele".

Em meados da década de 1950, haviam sido construídos modelos econométricos que se ajustavam bem aos dados de séries temporais, no sentido de que suas formas reduzidas (3) acompanhavam de perto os dados passados e se mostravam úteis na previsão de curto prazo. Além disso, por meio das restrições dos três tipos revisados acima, foi possível identificar seus parâmetros estruturais $A_j, B_j, R_k$. Usando essa estrutura estimada, foi possível simular os modelos para obter estimativas das consequências de diferentes políticas econômicas governamentais, como taxas de impostos, gastos ou política monetária.

Essa solução keynesiana para o problema de identificar um modelo estrutural tornou-se cada vez mais suspeita como resultado de desenvolvimentos de natureza tanto teórica quanto estatística. Muitos desses desenvolvimentos devem-se a esforços de pesquisadores simpáticos à tradição keynesiana, e muitos estavam bem avançados muito antes do fracasso espetacular dos modelos keynesianos na década de 1970.\footnote{Críticas às soluções keynesianas para o problema da identificação ao longo de linhas similares foram feitas por Lucas [17], Sims [33] e Sargent e Sims [31].}

Desde a sua criação, a macroeconomia tem sido criticada pela falta de "fundamentos em microeconomia e teoria do equilíbrio geral". Como observaram astutos comentaristas como Leontief [14] (com desaprovação) e Tobin [37] (com aprovação) desde cedo, a criação de um ramo distinto da teoria com seus próprios postulados distintos foi o objetivo consciente de Keynes. No entanto, um tema principal do trabalho teórico desde a \textit{Teoria Geral} tem sido a tentativa de usar a teoria microeconômica baseada no postulado clássico de que os agentes agem em seus próprios interesses para sugerir uma lista de variáveis que pertencem ao lado direito de um determinado cronograma comportamental, digamos, uma escala de demanda por um fator de produção ou uma escala de consumo.\footnote{[Esta nota foi adicionada na revisão, em parte como resposta aos comentários de Benjamin Friedman.] Muito deste trabalho foi feito por economistas operando bem dentro da tradição keynesiana, muitas vezes no contexto de algum modelo macroeconométrico keynesiano. Às vezes, recorreu-se a uma teoria com agentes otimizadores para resolver paradoxos empíricos, encontrando variáveis que haviam sido omitidas de algumas das formulações econométricas keynesianas anteriores. Os trabalhos de Modigliani e Friedman sobre consumo são bons exemplos dessa linha de trabalho, cujas implicações econométricas foram estendidas em trabalhos importantes por Robert Merton. Os trabalhos de Tobin e Baumol sobre equilíbrio de portfólio e de Jorgenson sobre investimento também estão na tradição de aplicar teorias microeconômicas de otimização para gerar relações de comportamento macroeconômico. Nos últimos trinta anos, os modelos econométricos keynesianos desenvolveram-se, em grande medida, na linha de tentar modelar o comportamento dos agentes como decorrente de problemas de otimização cada vez mais sofisticados. Nosso ponto aqui certamente não é afirmar que os economistas keynesianos tenham renunciado completamente a qualquer uso da teoria microeconômica otimizadora como guia. Pelo contrário, é que, especialmente quando problemas explicitamente estocásticos e dinâmicos foram estudados, tornou-se cada vez mais aparente que a teoria microeconômica tem implicações muito prejudiciais para as restrições convencionalmente usadas para identificar modelos macroeconométricos keynesianos. Além disso, como Tobin [37] enfatizou há muito tempo, há um ponto além do qual os modelos keynesianos devem suspender a hipótese de mercados em equilíbrio ou de agentes otimizadores se quiserem possuir as características operacionais e as implicações de política que são a marca registrada da economia keynesiana.} Mas, do ponto de vista da identificação de uma dada equação estrutural por meio de restrições do tipo (a), precisa-se de informação prévia confiável de que certas variáveis devem ser \textit{excluídas} do lado direito. A teoria microeconômica probabilística moderna quase nunca implica as restrições de exclusão que foram sugeridas por Keynes ou aquelas impostas pelos modelos macroeconométricos.

Para citar um exemplo que tem implicações extremamente graves para a identificação dos modelos macro existentes, as expectativas sobre os preços futuros, taxas de impostos e níveis de renda desempenham um papel crítico em muitos cronogramas de demanda e oferta nesses modelos. Por exemplo, nos melhores modelos, supõe-se tipicamente que a demanda por investimento responda às expectativas dos empresários sobre futuros créditos tributários, taxas de impostos e custos de fatores. Supõe-se tipicamente que a oferta de trabalho dependa da taxa de inflação que os trabalhadores esperam no futuro. Tais equações estruturais são geralmente identificadas pela premissa de que, por exemplo, a expectativa sobre o preço do fator ou a taxa de inflação atribuída aos agentes é uma função apenas de alguns valores defasados da própria variável que o agente supostamente está prevendo. No entanto, os próprios macromodelos contêm interações dinâmicas complicadas entre variáveis endógenas, incluindo os preços dos fatores e a taxa de inflação, e geralmente implicam que um agente sensato usaria valores atuais e muitos valores defasados de muitas e geralmente da maioria das variáveis endógenas e exógenas no modelo para formar expectativas sobre qualquer variável individual. Assim, virtualmente qualquer versão da hipótese de que os agentes se comportam de acordo com seus próprios interesses contradirá as restrições de identificação impostas à formação de expectativas. Além disso, as restrições sobre expectativas usadas para alcançar a identificação são inteiramente arbitrárias e não foram derivadas de nenhum princípio fundamental mais profundo sobre o comportamento econômico. Nenhum princípio geral jamais foi estabelecido que implicasse que, digamos, a taxa de inflação esperada devesse ser modelada como uma função linear apenas de taxas de inflação defasadas com pesos que somam a unidade; no entanto, essa hipótese é usada como uma restrição de identificação em quase todos os modelos existentes. O tratamento casual das expectativas não é um problema periférico nesses modelos, pois o papel das expectativas é onipresente neles e exerce uma influência maciça sobre suas propriedades dinâmicas (um ponto no qual o próprio Keynes insistia). O fracasso dos modelos existentes em derivar restrições sobre as expectativas a partir de quaisquer princípios básicos fundamentados na teoria econômica é um sintoma de um fracasso um pouco mais profundo e geral em derivar relações comportamentais de quaisquer problemas de otimização dinâmica consistentemente formulados.

Quanto à segunda categoria, restrições do tipo (b), os modelos macro keynesianos existentes fazem restrições severas \textit{a priori} sobre os $R_j$. Tipicamente, supõe-se que os $R_j$ sejam diagonais, de modo que a correlação serial defasada entre equações é ignorada e, além disso, assume-se que a ordem do processo $\epsilon_t$ seja curta, de modo que apenas a correlação serial de baixa ordem é permitida. Atualmente não existem fundamentos teóricos para introduzir essas restrições e, por boas razões, há pouca perspectiva de que a teoria econômica venha a fornecer tais fundamentos em breve. Em princípio, a identificação pode ser alcançada sem a imposição de tais restrições. Renunciar ao uso das restrições da categoria (b) aumentaria a necessidade de restrições das categorias (a) e (c). Em qualquer caso, os modelos macro existentes restringem pesadamente os $R$.

Voltando-nos para a terceira categoria, todos os grandes modelos existentes adotam uma classificação \textit{a priori} das variáveis nas categorias de estritamente endógenas, os $y_t$, e estritamente exógenas, os $x_t$. Cada vez mais, reconhece-se que a classificação de uma variável como "exógena" com base na observação de que ela \textit{poderia} ser definida sem referência aos valores atuais e passados de outras variáveis não tem nada a ver com a questão econometricamente relevante de como essa variável \textit{de fato} se relacionou com outras ao longo de um determinado período histórico. Além disso, à luz de desenvolvimentos recentes na econometria de séries temporais, sabemos que esse procedimento de classificação arbitrária não é necessário. Christopher Sims [34] mostrou que, em um contexto de séries temporais, a hipótese de exogeneidade econométrica pode ser testada. Ou seja, Sims mostrou que a hipótese de que $x_t$ é estritamente exógena econometricamente em (1) implica necessariamente certas restrições que podem ser testadas dadas séries temporais sobre os $y$ e os $x$. Testes seguindo a linha de Sims deveriam ser usados rotineiramente na verificação das categorizações em conjuntos de variáveis exógenas e endógenas. Até o momento, eles não têm sido. Construtores proeminentes de grandes modelos econométricos chegaram até a negar a utilidade de tais testes.\footnote{Por exemplo, veja o comentário de Albert Ando [35, especialmente pp. 209-210] e as observações de L. R. Klein [24].}

\section{Fracasso da Macroeconometria Keynesiana}

Nossa discussão na seção anterior levantou uma série de razões teóricas para acreditar que os parâmetros identificados como estruturais pelos métodos atualmente em uso na macroeconomia não são estruturais de fato. Ou seja, não há razão, em nossa opinião, para acreditar que esses modelos tenham isolado estruturas que permanecerão invariantes diante da classe de intervenções que figuram nas discussões contemporâneas sobre política econômica. No entanto, a questão de saber se um modelo específico é estrutural é empírica, não teórica. Se os modelos macroeconométricos tivessem compilado um histórico de estabilidade de parâmetros, particularmente diante de rupturas no comportamento estocástico das variáveis exógenas e distúrbios, seríamos céticos quanto à importância das objeções teóricas prévias do tipo que levantamos.

De fato, no entanto, o histórico dos principais modelos econométricos é, em qualquer dimensão que não seja a previsão incondicional de curtíssimo prazo, muito fraco. Testes estatísticos formais para instabilidade de parâmetros, conduzidos subdividindo séries passadas em períodos e verificando a estabilidade dos parâmetros ao longo do tempo, revelam invariavelmente grandes mudanças (para um exemplo, veja [23]). Além disso, essa dificuldade é implicitamente reconhecida pelos próprios construtores de modelos, que rotineiramente empregam um sistema elaborado de fatores de ajuste (\textit{add-factors}) na previsão, em uma tentativa de compensar a contínua "deriva" do modelo em relação às séries reais.

Embora não tenham sido projetados como tal por ninguém, os modelos macroeconométricos foram submetidos, na década de 1970, a um teste decisivo. Um elemento-chave em todos os modelos keynesianos é um "tradeoff" entre inflação e produto real: quanto maior a taxa de inflação, maior o produto (ou, equivalentemente, menor a taxa de desemprego). Por exemplo, os modelos do final da década de 1960 previam uma taxa de desemprego sustentada nos Estados Unidos de 4 por cento como consistente com uma taxa de inflação anual de 4 por cento. Muitos economistas na época instaram uma política deliberada de inflação com base nessa previsão. Certamente o caráter errático de "solavancos e paradas" da política real dos EUA na década de 1970 não pode ser atribuído a recomendações baseadas em modelos keynesianos, mas o viés inflacionário \textit{em média} da política monetária e fiscal nesse período deveria, de acordo com todos esses modelos, ter produzido as taxas médias de desemprego mais baixas de qualquer década desde a década de 1940. Na verdade, como sabemos, eles produziram o maior desemprego desde a década de 1930. Isso foi um fracasso econométrico em escala monumental.

Este fracasso não levou a conversões generalizadas de economistas keynesianos a outras crenças, nem se deveria esperar que levasse. Na economia, como em outras ciências, uma estrutura teórica é sempre mais ampla e flexível do que qualquer conjunto específico de equações, e há sempre a esperança de que, se um modelo específico falhar, pode-se encontrar um mais bem-sucedido baseado "grosseiramente" nas mesmas ideias. Isso teve, no entanto, já algumas consequências importantes, com implicações sérias tanto para a formulação de políticas econômicas quanto para a prática da ciência econômica.

Para a política econômica, o fato central é que as recomendações de política keynesianas não têm base mais sólida, em um sentido científico, do que as recomendações de economistas não keynesianos ou, aliás, de não economistas. Para notar uma consequência do amplo reconhecimento disso, a onda atual de sentimento protecionista dirigida a "salvar empregos" teria sido respondida, há dez anos, com o contra-argumento keynesiano de que a política fiscal pode atingir o mesmo fim, mas de forma mais eficiente. Hoje, é claro, ninguém levaria essa resposta a sério, por isso ela não é oferecida. De fato, economistas que há dez anos defendiam a política fiscal keynesiana como alternativa a controles diretos ineficientes favorecem cada vez mais estes últimos como "suplementos" à política keynesiana. A ideia parece ser que, se as pessoas se recusam a obedecer às equações que ajustamos ao seu comportamento passado, podemos aprovar leis para obrigá-las a fazê-lo.

Cientificamente, o fracasso keynesiano da década de 1970 resultou em uma nova abertura. Cada vez menos economistas estão envolvidos no monitoramento e refinamento dos grandes modelos econométricos; cada vez mais estão desenvolvendo teorias alternativas do ciclo econômico, baseadas em diferentes princípios teóricos. Além disso, maior atenção e respeito são concedidos às vítimas teóricas da Revolução Keynesiana, às ideias dos contemporâneos de Keynes e de economistas anteriores cujo pensamento foi considerado por anos como ultrapassado.

No momento atual, é impossível prever para onde esses desenvolvimentos levarão. Alguns, é claro, continuam a acreditar que os problemas dos modelos keynesianos existentes podem ser resolvidos dentro da estrutura atual, que esses modelos podem ser adequadamente refinados alterando algumas equações estruturais, adicionando ou subtraindo algumas variáveis aqui e ali, ou talvez desagregando vários blocos de equações. Formulamos nossas críticas precedentes em termos tão gerais precisamente para enfatizar seu caráter genérico e, consequentemente, a futilidade de buscar variações menores dentro dessa estrutura geral.

Uma segunda resposta ao fracasso dos métodos analíticos keynesianos é renunciar inteiramente aos métodos analíticos, retornando a métodos "baseados no julgamento". A primeira dessas respostas identifica os objetivos quantitativos e científicos da Revolução Keynesiana com os detalhes dos modelos específicos desenvolvidos até agora. A segunda renuncia tanto a esses modelos quanto aos objetivos que eles foram projetados para atingir. Existe, acreditamos, um curso intermediário, para o qual nos voltamos agora.

\section{Teoria do Ciclo Econômico de Equilíbrio}

Economistas anteriores à década de 1930 não reconheciam a necessidade de um ramo especial da economia, com seus próprios postulados especiais, projetado para explicar o ciclo econômico. Keynes fundou essa subdisciplina, chamada macroeconomia, porque pensou que era impossível explicar as características dos ciclos econômicos dentro da disciplina imposta pela teoria econômica clássica, uma disciplina imposta por sua insistência na adesão aos dois postulados: (a) que se assuma que os mercados se equilibram (\textit{clear}), e (b) que se assuma que os agentes agem em seu próprio interesse. O fato notável que parecia impossível reconciliar com esses dois postulados era a duração e a severidade das depressões econômicas e o desemprego em grande escala que elas acarretavam. Uma observação relacionada é que as medidas de demanda agregada e os preços estão positivamente correlacionados com as medidas de produto real e emprego, em aparente contradição com o resultado clássico de que mudanças em uma magnitude puramente nominal, como o nível geral de preços, eram meras "mudanças de unidade" que não deveriam alterar o comportamento real. Após libertar-se da "camisa de força" (ou disciplina) imposta pelos postulados clássicos, Keynes descreveu um modelo no qual regras práticas, como a função consumo e o cronograma de preferência pela liquidez, tomaram o lugar de funções de decisão que um economista clássico insistiria que fossem derivadas da teoria da escolha. E em vez de exigir que salários e preços fossem determinados pelo postulado de que os mercados se equilibram — o que para o mercado de trabalho parecia patentemente contradito pela severidade das depressões — Keynes tomou como um postulado não examinado que os salários nominais são "rígidos", significando que são definidos em um nível ou por um processo que poderia ser considerado não influenciado pelas forças macroeconômicas que ele propunha analisar.

Quando Keynes escreveu, os termos "equilíbrio" e "clássico" carregavam certas conotações positivas e normativas que pareciam excluir qualquer um dos modificadores de ser aplicado à teoria do ciclo econômico. O termo "equilíbrio" era pensado para referir-se a um sistema "em repouso", e tanto "equilíbrio" quanto "clássico" eram usados de forma intercambiável, por alguns, com "ideal". Assim, uma economia em equilíbrio clássico seria tanto imutável quanto não passível de melhorias por meio de intervenções políticas. Usando os termos dessa maneira, não é de admirar que poucos economistas considerassem a teoria do equilíbrio como um ponto de partida promissor para a compreensão dos ciclos econômicos e para o desenho de políticas para mitigá-los ou eliminá-los.

Nos últimos anos, o significado do termo "equilíbrio" passou por um desenvolvimento tão dramático que um teórico da década de 1930 não o reconheceria. É agora rotineiro descrever uma economia seguindo um processo estocástico multivariado como estando "em equilíbrio", o que significa nada mais do que em cada ponto do tempo os postulados (a) e (b) acima são satisfeitos. Este desenvolvimento, que decorreu principalmente do trabalho de K. J. Arrow [2] e G. Debreu [6], implica que simplesmente olhar para qualquer série temporal econômica e concluir que ela é um "fenômeno de desequilíbrio" é uma observação sem sentido. De fato, uma conjectura mais provável, com base no trabalho recente de Hugo Sonnenschein [36], é que a hipótese geral de que uma coleção de séries temporais descreve uma economia em equilíbrio competitivo é \textit{sem conteúdo}.\footnote{Para um exemplo que ilustra o vazio em um nível geral da afirmação de que "os empregadores estão sempre operando ao longo de demandas estocásticas dinâmicas por fatores", veja as observações sobre identificação econométrica em Sargent [29]. Em problemas aplicados que envolvem a modelagem das regras de decisão ótimas dos agentes, fica-se impressionado com o fato de que a generalização da especificação das funções objetivo dos agentes de formas plausíveis leva rapidamente à subidentificação econométrica. Um tipo de exemplo um pouco diferente é visto nas dificuldades de usar observações de séries temporais para refutar a visão de que "os agentes respondem apenas a mudanças inesperadas na oferta de moeda". Uma característica distintiva dos modelos macroeconométricos de equilíbrio descritos abaixo é que mudanças previsíveis na oferta monetária não afetam o PIB real ou o emprego total. Nos modelos keynesianos, mudanças previsíveis na oferta monetária de fato fazem com que o PIB real e o emprego se movam. Em um nível geral, é impossível discriminar entre essas duas visões observando séries temporais extraídas de uma economia descrita por um processo aleatório vetorial estacionário (Sargent [28]).}

A linha de pesquisa que está sendo perseguida por vários de nós envolve a tentativa de descobrir uma teoria de equilíbrio do ciclo econômico específica e econometricamente testável, que possa servir de base para a análise quantitativa da política macroeconômica. Não há como negar que essa abordagem é "contrarrevolucionária", pois pressupõe que Keynes e seus seguidores estavam errados ao desistir da possibilidade de que uma teoria de equilíbrio pudesse explicar o ciclo econômico. Até o momento, nenhum modelo macroeconométrico de equilíbrio bem-sucedido no nível de detalhe de, digamos, o modelo FMP, foi construído. Mas pequenos modelos teóricos de equilíbrio foram construídos, mostrando potencial para explicar algumas características-chave do ciclo econômico, há muito consideradas inexplicáveis dentro dos confins dos postulados clássicos. Os modelos de equilíbrio também fornecem razões para entender por que os modelos keynesianos estimados não se sustentam fora da amostra na qual foram estimados. Voltamo-nos agora para descrever alguns dos fatos fundamentais sobre os ciclos econômicos e a maneira como os novos modelos clássicos os enfrentam.

Por muito tempo, a maior parte da profissão econômica, com alguma razão, seguiu Keynes na rejeição dos modelos macroeconômicos clássicos porque eles pareciam incapazes de explicar algumas características importantes das séries temporais que medem agregados econômicos importantes. Talvez o fracasso mais importante do modelo clássico tenha sido sua aparente incapacidade de explicar a correlação positiva nas séries temporais entre preços e/ou salários, por um lado, e medidas de produto agregado ou emprego, por outro. Um segundo e relacionado fracasso foi sua incapacidade de explicar as correlações positivas entre medidas de demanda agregada, como o estoque monetário, e o produto agregado ou emprego. A análise estática dos modelos macroeconômicos clássicos tipicamente implicava que os níveis de produto e emprego eram determinados independentemente tanto do nível absoluto de preços quanto da demanda agregada. A presença pervasiva das correlações positivas mencionadas acima nas séries temporais parece consistente com conexões causais fluindo da demanda agregada e inflação para o produto e emprego, contrariamente às proposições de "neutralidade" clássica. Os modelos macroeconométricos keynesianos de fato implicam tais conexões causais.

Temos agora modelos teóricos rigorosos que ilustram como essas correlações podem emergir enquanto se mantêm os postulados clássicos de que os mercados se equilibram e os agentes otimizam.\footnote{Veja Edmund S. Phelps et al. [25] e Lucas [15], [16].} O passo fundamental na obtenção de tais modelos foi relaxar o postulado acessório usado em muita análise econômica clássica de que os agentes possuem informação perfeita. Os novos modelos clássicos continuam a assumir que os mercados sempre se equilibram e que os agentes otimizam. O postulado de que os agentes otimizam significa que suas decisões de oferta e demanda devem ser funções de variáveis reais, incluindo preços relativos percebidos. Assume-se que cada agente possui informação limitada e recebe informações sobre alguns preços com mais frequência do que sobre outros preços. Com base em sua informação limitada — as listas que possuem de preços absolutos atuais e passados de vários bens — assume-se que os agentes façam a melhor estimativa possível de todos os preços relativos que influenciam suas decisões de oferta e demanda. Como não possuem todas as informações que lhes permitiriam calcular perfeitamente os preços relativos com os quais se preocupam, os agentes cometem erros ao estimar os preços relativos pertinentes, erros que são inevitáveis dada a sua informação limitada. Em particular, sob certas condições, os agentes tenderão temporariamente a confundir um aumento geral em todos os preços absolutos como um aumento no preço relativo do bem que estão vendendo, levando-os a aumentar sua oferta daquele bem além do que haviam planejado anteriormente. Como todos estão, em média, cometendo o mesmo erro, o produto agregado subirá acima do que seria. Este aumento de produto acima do que seria ocorrerá sempre que o nível de preços médio de toda a economia deste período estiver acima do que os agentes esperavam que fosse com base na informação anterior. Simetricamente, o produto médio será reduzido sempre que o preço agregado se revelar inferior ao que os agentes esperavam. A hipótese de "expectativas racionais" está sendo imposta aqui porque se supõe que os agentes façam o melhor uso possível da informação limitada que possuem e assume-se que conheçam as distribuições de probabilidade objetivas pertinentes. Esta hipótese é imposta como forma de aderir aos princípios da teoria do equilíbrio.

Na teoria precedente, distúrbios na demanda agregada levam a uma correlação positiva entre mudanças inesperadas no nível de preços agregado e revisões no produto agregado em relação ao seu nível previamente planejado. Além disso, é um passo fácil mostrar que a teoria implica correlações entre revisões do produto agregado e mudanças inesperadas em quaisquer variáveis que ajudem a determinar a demanda agregada. Na maioria dos modelos macroeconômicos, a oferta monetária é um determinante da demanda agregada. A teoria precedente pode facilmente dar conta de correlações positivas entre revisões do produto agregado e aumentos inesperados na oferta monetária.

Embora tal teoria preveja correlações positivas entre a taxa de inflação ou a oferta monetária, por um lado, e o nível de produto, por outro, ela também afirma que essas correlações não representam "tradeoffs" que podem ser explorados por uma autoridade de política econômica. Ou seja, a teoria prevê que não há como a autoridade monetária seguir uma política ativista sistemática e atingir uma taxa de produto que seja, em média, mais alta ao longo do ciclo econômico do que a que ocorreria se ela simplesmente adotasse uma regra de $X$ por cento sem \textit{feedback}, do tipo que Friedman [8] e Simons [32] recomendaram. Pois a teoria prevê que o produto agregado é uma função de mudanças inesperadas atuais e passadas na oferta monetária. O produto será alto apenas quando a oferta monetária for e tiver sido maior do que se esperava que fosse, ou seja, maior que a média. Simplesmente não há como, em média, ao longo de todo o ciclo econômico, a oferta monetária ser maior que a média. Assim, embora a teoria precedente seja capaz de explicar algumas das correlações que há muito se pensava invalidarem a teoria macroeconômica clássica, a teoria é clássica tanto em sua adesão aos postulados teóricos clássicos quanto no sabor "não ativista" de suas implicações para a política monetária.

Modelos econométricos de pequena escala no sentido da Seção 2 deste artigo foram construídos capturando algumas das principais características dos modelos de equilíbrio descritos acima.\footnote{Por exemplo, Sargent [27]. A insatisfação com os métodos keynesianos de alcançar a identificação também levou a outras linhas de trabalho macroeconométrico. Uma linha são os "modelos de índices" descritos por Sargent e Sims [31] e Geweke [10]. Estes modelos representam uma forma estatisticamente precisa de implementar a noção de Wesley Mitchell de que existe um pequeno número de influências comuns que explicam a covariação de um grande número de agregados econômicos ao longo do ciclo econômico. Esta hipótese de "baixa dimensionalidade" é um dispositivo potencial para restringir o número de parâmetros a serem estimados em modelos de séries temporais vetoriais. Esta linha de trabalho não é inteiramente ateórica (mas veja os comentários de Ando e Klein em Sims [35]), embora seja distintamente não keynesiana. Por coincidência, certos modelos de equilíbrio do ciclo econômico parecem levar a modelos de índices de baixa dimensionalidade com um padrão interessante de cargas das variáveis nos índices. Em geral, os modelos keynesianos modernos não assumem tão facilmente uma forma de índice baixo. Veja a discussão em Sargent e Sims [31].} Em particular, esses modelos incorporam a hipótese de que as expectativas são racionais, ou que toda a informação disponível é utilizada pelos agentes. Até certo ponto, esses modelos alcançam a identificação econométrica invocando restrições em cada uma das três categorias (a), (b) e (c). No entanto, uma característica distintiva destes modelos "clássicos" é que eles também dependem fortemente de uma importante quarta categoria de restrições de identificação. Esta categoria (d) consiste em um conjunto de restrições que são derivadas da teoria econômica probabilística, mas não desempenham papel algum na estrutura keynesiana. Estas restrições em geral não assumem a forma de restrições de zero do tipo (a). Em vez disso, as restrições da teoria assumem tipicamente a forma de restrições entre equações entre os parâmetros $A_j, B_j, C_j$. A fonte dessas restrições é a implicação da teoria econômica de que as decisões atuais dependem das previsões dos agentes sobre variáveis futuras, combinada com a implicação de que essas previsões são formadas de maneira ótima, dado o comportamento das variáveis passadas. Essas restrições não possuem uma expressão matemática tão simples quanto meramente definir um número de parâmetros igual a zero, mas sua motivação econômica é fácil de entender. Formas de utilizar essas restrições na estimação e teste econométrico estão sendo rapidamente desenvolvidas.

Outra característica fundamental do trabalho recente em modelos macroeconométricos de equilíbrio é que a dependência de categorizações inteiramente \textit{a priori} (c) de variáveis como estritamente exógenas e endógenas foi marcadamente reduzida, embora não totalmente eliminada. Este desenvolvimento decorre conjuntamente do fato de que os modelos atribuem papéis importantes às previsões ótimas dos agentes sobre variáveis futuras e da demonstração de Christopher Sims de que existe uma conexão estreita entre o conceito de exogeneidade econométrica estrita e as formas dos preditores ótimos para um vetor de séries temporais. Construir um modelo com expectativas racionais força necessariamente a considerar qual conjunto de outras variáveis ajuda a prever uma determinada variável, digamos, a renda ou a taxa de inflação. Se a variável $y$ ajuda a prever a variável $x$, então os teoremas de Sims implicam que $x$ não pode ser considerada exógena em relação a $y$. O resultado desta conexão entre previsibilidade e exogeneidade foi que, nos modelos macroeconométricos de equilíbrio, a distinção entre variáveis endógenas e exógenas não foi traçada em uma base inteiramente \textit{a priori}. Além disso, casos especiais dos modelos teóricos, que frequentemente envolvem restrições laterais sobre os $R_j$ não derivadas da teoria econômica, possuem fortes predições testáveis quanto às relações de exogeneidade entre as variáveis.

Uma característica fundamental dos modelos macroeconométricos de equilíbrio é que, como resultado das restrições entre as matrizes $A_j, B_j$ e $C_j$, os modelos preveem que, em geral, os parâmetros em muitas das equações mudarão se houver uma intervenção de política que assuma a forma de uma mudança em uma equação que descreve como alguma variável de política está sendo definida. Como ignoram essas restrições entre equações, os modelos keynesianos em geral assumem que todas as outras equações permanecem inalteradas quando uma equação que descreve uma variável de política é alterada. Nossa visão é que esta é uma razão importante pela qual os modelos keynesianos quebraram quando ocorreram mudanças importantes nas equações que governam as variáveis de política ou variáveis exógenas. Nossa esperança é que os métodos que descrevemos nos deem a capacidade de prever as consequências para todas as equações de mudanças nas regras que governam as variáveis de política. Ter essa capacidade é necessário antes que possamos alegar ter uma base científica para fazer afirmações quantitativas sobre a política macroeconômica.

Atualmente, esses novos desenvolvimentos teóricos e econométricos não foram totalmente integrados, embora esteja claro que estão muito próximos, tanto conceitual quanto operacionalmente. Nossa preferência seria considerar os melhores modelos de equilíbrio existentes atualmente como protótipos de modelos melhores e futuros que, esperamos, se mostrarão de utilidade prática na formulação de políticas. Mas não devemos subestimar o sucesso econométrico já alcançado pelos modelos de equilíbrio. Versões iniciais desses modelos foram estimadas e submetidas a alguns testes econométricos rigorosos por McCallum [20], Barro [3], [4] e Sargent [27], com o resultado de que parecem de fato capazes de explicar algumas características amplas do ciclo econômico. Modelos novos e mais sofisticados envolvendo restrições entre equações mais complicadas estão em andamento (Sargent [29]). O trabalho realizado até o momento já mostrou que os modelos de equilíbrio são capazes de atingir ajustes dentro da amostra quase tão bons quanto os obtidos pelos modelos keynesianos, tornando assim concreto o ponto de que os bons ajustes dos modelos keynesianos não fornecem nenhuma boa razão para confiar nas recomendações de política delas derivadas.

\section{Crítica à Teoria de Equilíbrio}

A ideia central das explicações de equilíbrio dos ciclos econômicos, conforme esboçada acima, é que as flutuações econômicas surgem à medida que os agentes reagem a mudanças não antecipadas em variáveis que afetam suas decisões. Está claro que qualquer explicação deste tipo geral deve trazer consigo limitações severas na capacidade da política governamental de compensar estas mudanças iniciais. Primeiro, os governos devem, de alguma forma, ter a capacidade de prever choques que são invisíveis para os agentes privados, mas ao mesmo tempo não ter a capacidade de revelar esta informação antecipada (evitando assim os choques). Embora não seja difícil escrever modelos teóricos nos quais estas duas condições são assumidas como válidas, é difícil imaginar situações reais nas quais tais modelos se aplicariam. Segundo, a política anticíclica governamental deve ser, ela própria, não antecipável pelos agentes privados (certamente uma condição frequentemente realizada historicamente), enquanto, ao mesmo tempo, estar sistematicamente relacionada ao estado da economia. A eficácia reside, então, na incapacidade dos agentes privados de reconhecer padrões sistemáticos na política monetária e fiscal.

Em grande medida, a crítica aos modelos de equilíbrio é simplesmente uma reação a estas implicações para a política econômica. Tão amplo é (ou era) o consenso de que a tarefa da macroeconomia é a descoberta das políticas monetárias e fiscais específicas que podem eliminar as flutuações reagindo à instabilidade do setor privado que a afirmação de que esta tarefa não deve ou não pode ser realizada é considerada frívola, independentemente de qualquer raciocínio e evidência que a sustente. Certamente deve-se ter alguma simpatia por esta reação: uma fé infundada na cura de um mal específico serviu muitas vezes de estímulo para a descoberta de curas genuínas. No entanto, confundir uma fé possivelmente funcional na existência de políticas monetárias e fiscais reativas e eficazes com evidência científica de que tais políticas são conhecidas é claramente perigoso, e usar tal fé como critério para julgar a extensão em que teorias específicas "se ajustam aos fatos" é ainda pior.

Existem, é claro, questões legítimas envolvendo a capacidade das teorias de equilíbrio de se ajustarem aos fatos do ciclo econômico. De fato, esta é a razão de nossa insistência no caráter preliminar e tentativo dos modelos específicos que temos agora. No entanto, estes modelos tentativos compartilham certas características que podem ser consideradas essenciais, de modo que não é irracional especular sobre a probabilidade de que qualquer modelo deste tipo possa ser bem-sucedido, ou perguntar: o que os teóricos do ciclo econômico de equilíbrio terão daqui a dez anos se tivermos sorte?

Quatro razões gerais para pessimismo que têm sido proeminentemente avançadas são: (a) o fato de que os modelos de equilíbrio postulam mercados equilibrados; (b) a afirmação de que estes modelos não podem explicar a "persistência" (correlação serial) dos movimentos cíclicos; (c) o fato de que os modelos implementados econometricamente são lineares (em logaritmos); e (d) o fato de que o comportamento de aprendizado não foi incorporado. Discutimos cada um destes pontos em subseções distintas.

\subsection{Mercados em Equilíbrio}

Uma característica essencial dos modelos de equilíbrio é que todos os mercados se equilibram (\textit{clear}), ou que todos os preços e quantidades observados sejam explicáveis como resultados de decisões tomadas por firmas e famílias individuais. Na prática, isso significou uma suposição convencional e competitiva de oferta igual à demanda, embora outros tipos de equilíbrio possam ser facilmente imaginados (se não tão facilmente analisados). Se, portanto, alguém tomar como um "fato" básico que os mercados de trabalho não se equilibram, chega-se imediatamente a uma contradição entre teoria e fato. Os fatos que realmente temos, no entanto, são simplesmente as séries temporais disponíveis sobre emprego e taxas salariais, além das respostas às nossas pesquisas de desemprego. Mercados equilibrados é apenas um princípio, não verificável por observação direta, que pode ou não ser útil na construção de hipóteses bem-sucedidas sobre o comportamento dessas séries. Princípios alternativos, como o postulado da existência de um leiloeiro terceiro induzindo "rigidez" salarial e mercados não equilibrados, são similarmente "irrealistas", no sentido não especialmente importante de não oferecer uma boa descrição das instituições observadas do mercado de trabalho.

Um refinamento do postulado não explicado de um mercado de trabalho não equilibrado foi sugerido pelo fato indiscutível de que existem contratos de trabalho de longo prazo com horizontes de dois ou três anos. No entanto, a duração \textit{per se} ao longo da qual os contratos ocorrem não afeta a questão, pois sabemos, de Arrow e Debreu, que se contratos infinitamente longos forem determinados de modo que preços e salários dependam da mesma informação que está disponível sob a suposição de equilíbrio de mercado período a período, então resultará precisamente o mesmo processo de preço-quantidade com o contrato de longo prazo que ocorreria sob equilíbrio de mercado período a período. Assim, a teorização de equilíbrio fornece uma maneira, provavelmente a única que temos, de construir um modelo de um contrato de longo prazo. O fato de existirem contratos de longo prazo, então, não tem implicações sobre a aplicabilidade da teorização de equilíbrio. Pelo contrário, a questão real aqui é se os contratos reais podem ser adequadamente explicados dentro de um modelo de equilíbrio, ou seja, um modelo no qual os agentes agem em seus próprios interesses. Stanley Fischer [7], Edmund Phelps e John Taylor [26] e Robert Hall [12] mostraram que algumas das conclusões "não ativistas" dos modelos de equilíbrio são modificadas se for substituído o equilíbrio de mercado período a período pela imposição de contratos de longo prazo baseados em conjuntos de informações restritos que são impostos exogenamente e que se assume serem independentes dos regimes monetários e fiscais. A teoria econômica nos leva a prever que os custos de coleta e processamento de informações tornarão ótimo que os contratos dependam de um pequeno subconjunto de informações que poderiam ser coletadas em qualquer data. Mas a teoria também sugere que o conjunto particular de informações sobre o qual os contratos serão baseados não é imutável, mas depende da estrutura de custos e benefícios para a coleta de vários tipos de informações. Esta estrutura de custos e benefícios mudará a cada mudança nos processos estocásticos exógenos enfrentados pelos agentes. Esta presunção teórica é apoiada por um exame da maneira como os contratos de trabalho diferem entre países de alta e baixa inflação e a maneira como evoluíram nos Estados Unidos nos últimos 25 anos.

Portanto, a questão aqui é realmente a mesma fundamental envolvida na disputa entre Keynes e os economistas clássicos: é adequado considerar certas características superficiais dos contratos salariais existentes como dadas ao analisar as consequências de regimes monetários e fiscais alternativos? A teoria econômica clássica nega que essas características possam ser tomadas como dadas. Para entender as implicações dos contratos de longo prazo para a política monetária, precisa-se de um modelo da maneira como esses contratos provavelmente responderão a regimes de política monetária alternativos. Uma extensão dos modelos de equilíbrio existentes nesta direção pode muito bem levar a variações interessantes, mas parece-nos improvável que modificações importantes nas implicações destes modelos para a política monetária e fiscal decorram disso.

\subsection{Persistência}

Uma segunda linha de crítica decorre da observação correta de que, se as expectativas dos agentes forem racionais e se seus conjuntos de informações incluírem valores defasados da variável que está sendo prevista, então os erros de previsão dos agentes devem ser um processo aleatório serialmente não correlacionado. Ou seja, em média, não deve haver relações detectáveis entre o erro de previsão deste período e qualquer erro de previsão de período anterior. Esta característica levou vários críticos a concluir que os modelos de equilíbrio são incapazes de explicar mais do que uma parte insignificante dos movimentos altamente correlacionados serialmente que observamos no produto real, emprego, desemprego e outras séries. Tobin colocou o argumento de forma sucinta em [38]:

\begin{quote}
    Uma explicação atualmente popular para variações no emprego é a confusão temporária de preços relativos e absolutos. Empregadores e trabalhadores são enganados por muitos empregos devido à inflação inesperada, mas apenas até que aprendam que ela afeta outros preços, e não apenas os preços do que vendem. O inverso acontece temporariamente quando a inflação fica abaixo da expectativa. Este modelo mal pode explicar mais do que um desequilíbrio transitório nos mercados de trabalho.
    
    Então, como podem os fiéis explicar os ciclos lentos de desemprego que realmente observamos? Apenas argumentando que a própria taxa natural flutua, que as variações nas taxas de desemprego são substancialmente mudanças no desemprego voluntário, friccional ou estrutural, em vez de desemprego involuntário devido à demanda geralmente deficiente.
\end{quote}

Os críticos tipicamente concluem que a teoria atribui apenas um papel muito menor às flutuações da demanda agregada e depende necessariamente de distúrbios na oferta agregada para explicar a maior parte das flutuações no produto real ao longo do ciclo econômico. Como Modigliani [21] caracterizou as implicações da teoria: "Em outras palavras, o que aconteceu com os Estados Unidos na década de 1930 foi um ataque severo de preguiça contagiosa".

Esta crítica é falaciosa porque não distingue adequadamente entre "fontes de impulsos" e "mecanismos de propagação", uma distinção enfatizada por Ragnar Frisch em um artigo clássico de 1933 [9] que forneceu muitos dos fundamentos técnicos para os modelos macroeconométricos keynesianos. Embora a nova teoria clássica implique que os erros de previsão, que são os "impulsos" da demanda agregada, sejam serialmente não correlacionados, é certamente logicamente possível que existam "mecanismos de propagação" em ação que convertam esses impulsos em movimentos serialmente correlacionados em variáveis reais como produto e emprego. De fato, dois mecanismos de propagação concretos já foram mostrados em trabalhos teóricos detalhados como sendo capazes de desempenhar precisamente essa função. Um mecanismo decorre da presença de custos para as firmas de ajustar rapidamente seus estoques de capital e trabalho. A presença desses custos é conhecida por tornar ótimo para as firmas espalhar ao longo do tempo sua resposta aos sinais de preços relativos que recebem. No contexto atual, tal mecanismo faz com que uma firma converta os erros de previsão serialmente não correlacionados na predição de preços relativos em movimentos serialmente correlacionados nas demandas de fatores e no produto.

Um segundo mecanismo de propagação já está presente na maioria dos modelos clássicos de crescimento econômico. Sabe-se que os planos ótimos de acumulação das famílias para direitos sobre capital físico e outros ativos converterão impulsos serialmente não correlacionados em demandas serialmente correlacionadas para a acumulação de ativos reais. Isso acontece porque os agentes tipicamente desejarão dividir quaisquer mudanças inesperadas nos preços ou renda enfrentados pelos agentes. Essa dependência faz com que o estoque de ativos do próximo período dependa dos estoques iniciais e das mudanças inesperadas nos preços ou na renda enfrentados pelos agentes. Essa dependência faz com que surpresas serialmente não correlacionadas levem a movimentos serialmente correlacionados nas demandas por ativos físicos. Lucas [16] mostrou como este mecanismo de propagação aceita prontamente erros na previsão da demanda agregada como uma fonte de "impulso".

Um terceiro provável mecanismo de propagação é identificado pelo trabalho recente em teoria da busca.\footnote{Por exemplo [19], [22] e [18].} A teoria da busca fornece uma explicação para o motivo pelo qual trabalhadores que, por algum motivo, se veem sem emprego, acharão racional não necessariamente aceitar a primeira oferta de emprego que aparecer, mas permanecer desempregados por algum período até que uma oferta melhor se materialize. Da mesma forma, a teoria fornece razões pelas quais uma firma pode achar ótimo esperar até que um candidato a emprego mais adequado apareça, de modo que as vagas persistam por algum tempo. Ao contrário dos dois primeiros mecanismos de propagação mencionados, modelos teóricos consistentes que permitam que este mecanismo aceite erros na previsão da demanda agregada como um impulso ainda não foram totalmente elaborados por razões principalmente técnicas, mas parece provável que este mecanismo venha a desempenhar um papel importante em um modelo bem-sucedido do comportamento das séries temporais da taxa de desemprego.

Em modelos onde os agentes possuem informação imperfeita, qualquer um dos dois primeiros e muito provavelmente o terceiro mecanismo é capaz de fazer com que movimentos serialmente correlacionados em variáveis reais decorram da introdução de uma sequência de erros de previsão serialmente não correlacionados. Assim, modelos teóricos e econométricos foram construídos nos quais, em princípio, o processo de erros de previsão serialmente não correlacionados é capaz de responder por qualquer proporção entre zero e um da variância de estado estacionário do produto real ou do emprego. O argumento de que tais modelos devem necessariamente atribuir a maior parte da variância no produto real e emprego a variações na oferta agregada é simplesmente logicamente errado.

\subsection{Linearidade}

A maior parte do trabalho econométrico que implementa modelos de equilíbrio envolveu o ajuste de modelos estatísticos lineares nas variáveis (mas frequentemente altamente não lineares nos parâmetros). Esta característica é passível de crítica com base no princípio indiscutível de que existem geralmente modelos não lineares que fornecem melhores aproximações do que os modelos lineares. Mais especificamente, modelos lineares nas variáveis não fornecem método algum para detectar e analisar efeitos sistemáticos de momentos de ordem superior aos primeiros choques e variáveis exógenas sobre os primeiros momentos das variáveis endógenas. Tais efeitos sistemáticos estão geralmente presentes onde as variáveis endógenas são definidas por agentes avessos ao risco.

Não há razão \textit{teórica} para que a maior parte do trabalho aplicado tenha usado modelos lineares, apenas razões técnicas convincentes dada a tecnologia de computação atual. O requisito técnico predominante do trabalho econométrico que impõe expectativas racionais é a capacidade de escrever expressões analíticas que forneçam as regras de decisão dos agentes como funções dos parâmetros de suas funções objetivo e como funções dos parâmetros que regem os processos aleatórios exógenos que eles enfrentam. Problemas dinâmicos de máximo estocástico com objetivos quadráticos, que dão origem a regras de decisão lineares, atendem a esse requisito essencial, o que é sua virtude. Apenas algumas outras formas funcionais para as funções objetivo dos agentes em problemas dinâmicos de ótimo estocástico têm essa mesma tratabilidade analítica necessária. A tecnologia de computação no futuro previsível parece exigir o trabalho com tal classe de funções, e a classe de regras de decisão lineares pareceu apenas a mais conveniente para a maioria dos propósitos. Nenhuma questão de princípio está envolvida na seleção de uma da classe muito restrita de funções disponíveis para nós. Teoricamente, sabemos como calcular por meio de métodos recursivos dispendiosos as regras de decisão não lineares que decorreriam de uma classe muito ampla de funções objetivo; nenhum novo princípio econométrico estaria envolvido na estimativa de seus parâmetros, apenas uma conta de computador muito mais alta. Além disso, como Frisch e Slutsky enfatizaram, equações de diferença estocásticas lineares parecem um dispositivo muito flexível para estudar ciclos econômicos. É uma questão em aberto se para explicar as características centrais do ciclo econômico haverá uma grande recompensa no ajuste de modelos não lineares.

\subsection{Modelos Estacionários e a Negligência do Aprendizado}

Benjamin Friedman e outros criticaram os modelos de expectativas racionais aparentemente sob o argumento de que grande parte do trabalho teórico e quase todo o trabalho empírico assumiram que os agentes têm operado por um longo tempo em um ambiente estocasticamente estacionário. Como consequência, tipicamente assume-se que os agentes descobriram as leis de probabilidade das variáveis que desejam prever. Como Modigliani apresentou o argumento em [21]:

\begin{quote}
    Ao nível lógico, Benjamin Friedman chamou a atenção para a omissão [dos modelos macroeconômicos de equilíbrio] de um mecanismo explícito de aprendizado, e sugeriu que, como resultado, isso só pode ser interpretado como uma descrição não de equilíbrio de curto prazo, mas de longo prazo, no qual nenhum agente desejaria recontratar. Mas então as implicações [dos modelos macroeconômicos de equilíbrio] são claramente nada surpreendentes, e sua relevância política é quase nula (p. 6).
\end{quote}

Mas tem sido apenas uma questão de conveniência analítica e não de necessidade que os modelos de equilíbrio tenham usado a suposição de "choques" stocasticamente estacionários e a suposição de que os agentes já aprenderam as distribuições de probabilidade que enfrentam. Ambas as suposições podem ser abandonadas, embora a um custo em termos da simplicidade do modelo.\footnote{Por exemplo, veja Crawford [5] e Grossman [11].} De fato, dentro da estrutura de funções objetivo quadráticas, na qual se aplica o "princípio da separação", pode-se aplicar a "fórmula de filtragem de Kalman" para derivar regras de decisão lineares ótimas com coeficientes dependentes do tempo. Nesta estrutura, o "filtro de Kalman" permite uma aplicação elegante do aprendizado bayesiano para atualizar as regras de previsão ótimas de período a período, à medida que novas informações se tornam disponíveis. O filtro de Kalman também permite a derivação de regras de decisão ótimas para uma classe interessante de processos exógenos não estacionários assumidos como enfrentados pelos agentes. A teorização de equilíbrio neste contexto conduz, assim, prontamente a um modelo de como a não estacionariedade do processo e o aprendizado bayesiano aplicado pelos agentes às variáveis exógenas levam a coeficientes dependentes do tempo nas regras de decisão dos agentes.

Embora modelos que incorporam o aprendizado bayesiano e a não estacionariedade estocástica sejam tecnicamente viáveis e consistentes com a estratégia de modelagem de equilíbrio, quase nenhum trabalho aplicado bem-sucedido nessas linhas veio à luz. Uma razão é provavelmente que modelos de séries temporais não estacionárias são pesados e existem em muitas variedades. Outra é que a hipótese do aprendizado bayesiano é vazia até que se impute arbitrariamente uma distribuição prévia (\textit{prior}) aos agentes ou se desenvolva um método de estimar os parâmetros da prévia a partir de dados de séries temporais. Determinar uma distribuição prévia a partir dos dados envolveria a estimativa de uma série de condições iniciais e proliferaria parâmetros de perturbação de uma maneira muito desagradável. É uma questão empírica se estas técnicas valerão a pena em termos de explicação de séries temporais macroeconômicas; não é uma questão de distinção entre o equilíbrio e os modelos macroeconométricos keynesianos. Na verdade, nenhum modelo macroeconométrico keynesiano existente incorpora um modelo econômico de aprendizado ou um modelo econômico que restrinja de qualquer maneira o padrão de não estacionariedades de coeficientes entre as equações.

Os modelos macroeconométricos criticados por Friedman e Modigliani, que assumem que os agentes "pegaram o jeito" dos processos aleatórios estacionários que enfrentam, dão origem a sistemas de equações de diferença estocásticas lineares da forma (1), (2) e (4). Como se sabe há muito tempo, tais equações de diferença estocásticas geram séries que "se parecem" com séries temporais econômicas. Além disso, se visualizados como estruturais (ou seja, invariantes em relação às intervenções de política), os modelos têm algumas das implicações para a política anticíclica que descrevemos acima. Se estas implicações de política estão ou não corretas depende de se os modelos são estruturais ou não, e de forma alguma depende de se os modelos podem ser caricaturados com termos como "longo prazo" ou "curto prazo".

Vale a pena reenfatizar que não desejamos que nossas respostas a estas críticas sejam confundidas com uma alegação de que os modelos de equilíbrio existentes podem explicar satisfatoriamente todas as principais características do ciclo econômico observado. Pelo contrário, argumentamos que nenhuma razão sólida foi ainda avançada que sequer sugira que estes modelos sejam, como classe, incapazes de fornecer uma teoria satisfatória do ciclo econômico.

\section{Resumo e Conclusões}

Tentemos expor de forma compacta os principais argumentos avançados neste artigo. Em seguida, comentaremos brevemente as principais implicações destes argumentos para a maneira como podemos pensar utilmente sobre a política econômica.

Primeiro, e mais importante, os modelos macroeconométricos keynesianos existentes são incapazes de fornecer orientação confiável na formulação de políticas monetárias, fiscais e outras. Esta conclusão baseia-se em parte nos espetaculares fracassos recentes destes modelos e, em parte, na sua falta de uma base teórica ou econométrica sólida. Segundo, neste último terreno, não há esperança de que modificações menores ou mesmo maiores destes modelos levem a uma melhoria significativa na sua confiabilidade.

Terceiro, modelos de equilíbrio podem ser formulados de forma livre destas dificuldades e oferecer um conjunto diferente de princípios que podem ser usados para identificar modelos econométricos estruturais. Os elementos-chave destes modelos são que os agentes são racionais, reagindo a mudanças de política de uma maneira que atenda aos seus melhores interesses privados, e que os impulsos que desencadeiam as flutuações econômicas são principalmente choques não antecipados.

Quarto, os modelos de equilíbrio já desenvolvidos explicam as principais características qualitativas do ciclo econômico. Estes modelos estão sendo submetidos a críticas contínuas, especialmente por aqueles envolvidos no seu desenvolvimento, mas os argumentos no sentido de que as teorias de equilíbrio são, em princípio, incapazes de explicar uma parte substancial das flutuações observadas parecem dever-se principalmente a simples mal-entendidos.

As implicações de política das teorias de equilíbrio são por vezes caricaturadas, por comentadores amigos e hostis, como a afirmação de que "a política econômica não importa" ou "não tem efeito".\footnote{Uma fonte principal desta crença é provavelmente Sargent e Wallace [30], no qual se mostrou que, no contexto de um modelo macroeconômico bastante padrão, mas com expectativas dos agentes assumidas racionais, a escolha de uma regra monetária reativa não tem consequência para o comportamento das variáveis reais. O ponto deste exemplo era mostrar que, precisamente dentro daquele modelo usado para racionalizar políticas monetárias reativas, tais políticas podiam ser mostradas como sem valor. Dificilmente se segue que toda política seja ineficaz em todos os contextos.} Esta implicação certamente assustaria os economistas neoclássicos que aplicaram com sucesso a teoria do equilíbrio ao estudo de inúmeros problemas envolvendo efeitos importantes das políticas fiscais na alocação de recursos e distribuição de renda. Nossa intenção não é rejeitar estas realizações, mas sim tentar \textit{imitá-las}, ou estender os métodos de equilíbrio que foram aplicados a muitos problemas econômicos para cobrir um fenômeno que até agora resistiu à sua aplicação: o ciclo econômico.

Caso esta arbitragem intelectual se prove bem-sucedida, sugerirá mudanças importantes na maneira como pensamos sobre a política econômica. Mais fundamentalmente, ela direciona a atenção para a necessidade de pensar na política como a escolha de "regras do jogo" estáveis, bem compreendidas pelos agentes econômicos. Apenas em tal ambiente a teoria econômica nos ajudará a prever as ações que os agentes escolherão tomar. Segundo, esta abordagem sugere que as políticas que afetam o comportamento principalmente porque suas consequências não podem ser corretamente diagnosticadas, como a instabilidade monetária e o financiamento do déficit, têm apenas a capacidade de perturbar. A provisão deliberada de desinformação não pode ser usada de maneira sistemática para melhorar o ambiente econômico.

Os objetivos da teoria do ciclo econômico de equilíbrio são tomados, sem modificação, do objetivo que motivou a construção dos modelos macroeconométricos keynesianos: fornecer um meio cientificamente fundamentado de avaliar, quantitativamente, os prováveis efeitos de políticas econômicas alternativas. Sem os sucessos econométricos alcançados pelos modelos keynesianos, este objetivo seria simplesmente inconcebível. A menos que os limites agora evidentes destes modelos também sejam francamente reconhecidos e direções radicalmente novas sejam tomadas, as realizações reais da Revolução Keynesiana serão perdidas tão certamente quanto aquelas que agora sabemos serem ilusórias.

\section*{REFERÊNCIAS}

\begin{enumerate}
    \item[{[1]}] Ando, Albert. "A Comment," em [35].
    \item[{[2]}] Arrow, Kenneth J. "The Role of Securities in the Optimal Allocation of Risk Bearing." \textit{Review of Economic Studies} 31 (1964): 91-96.
    \item[{[3]}] Barro, Robert J. "Unanticipated Money Growth and Unemployment in the United States." \textit{American Economic Review} 67 (1977): 101-15.
    \item[{[4]}] \underline{\hspace{1cm}}. "Unanticipated Money, Output and the Price Level in the United States." \textit{Journal of Political Economy} (no prelo).
    \item[{[5]}] Crawford, Robert. "Implications of Learning for Economic Models of Uncertainty." Manuscrito. Pittsburgh: Carnegie-Mellon University, 1971.
    \item[{[6]}] Debreu, Gerard. \textit{The Theory of Value}. New York: Wiley, 1959.
    \item[{[7]}] Fischer, Stanley. "Long-term Contracts, Rational Expectations, and the Optimal Money Supply Rule." \textit{Journal of Political Economy} 85 (1977): 191-206.
    \item[{[8]}] Friedman, Milton. "A Monetary and Fiscal Framework for Economic Stability." \textit{American Economic Review} 38 (1948): 245-64.
    \item[{[9]}] Frisch, Ragnar. "Propagation Problems and Impulse Problems in Dynamic Economics." Republicado em \textit{AEA Readings in Business Cycles}, editado por R.A. Gordon e L.R. Klein. Vol. X, 1965.
    \item[{[10]}] Geweke, John. "The Dynamic Factor Analysis of Economic Time Series." Em \textit{Latent Variables In Socio-Economic Models}, editado por D. Aigner e A. Goldberger, pp. 365-383. Amsterdam: North Holland, 1977.
    \item[{[11]}] Grossman, Sanford. "Rational Expectations and the Econometric Modeling of Markets Subject to Uncertainty: A Bayesian Approach." \textit{Journal of Econometrics} 3 (1975): 255-272.
    \item[{[12]}] Hall, Robert E. "The Macroeconomic Impact of Changes in Income Taxes in the Short and Medium Runs." \textit{Journal of Political Economy} 86 (1978): S71-S86.
    \item[{[13]}] Keynes, J. M. \textit{The General Theory of Employment, Interest, and Money}. London: Macmillan, 1936.
    \item[{[14]}] Leontief, W. "Postulates: Keynes' General Theory and the Classicists." Em \textit{The New Economics, Keynes' Influences on Theory and Public Policy}, editado por S. Harris. Clifton, New Jersey: Augustus Kelley, 1965.
    \item[{[15]}] Lucas, R. E., Jr. "Expectations and the Neutrality of Money." \textit{Journal of Economic Theory} 4, No. 2 (Abril 1972): 102-123.
    \item[{[16]}] \underline{\hspace{1cm}}. "An Equilibrium Model of the Business Cycle." \textit{Journal of Political Economy} 83, No. 6 (Dezembro 1975): 1113-1144.
    \item[{[17]}] \underline{\hspace{1cm}}. "Econometric Policy Evaluation: A Critique." Em \textit{The Phillips Curve and Labor Markets}, editado por K. Brunner e A. H. Meltzer. Carnegie-Rochester Conference Series on Public Policy, 1: 19-46. Amsterdam: North Holland, 1976.
    \item[{[18]}] \underline{\hspace{1cm}} e Prescott, Edward C. "Equilibrium Search and Unemployment." \textit{Journal of Economic Theory} 7 (1974): 188-209.
    \item[{[19]}] McCall, John. "The Economics of Information and Optimal Stopping Rules." \textit{Journal of Business} 38 (1965): 300-317.
    \item[{[20]}] McCallum, Bennett. "Rational Expectations and the Natural Rate Hypothesis: Some Consistent Estimates." \textit{Econometrica} 44 (1976): 43-52.
    \item[{[21]}] Modigliani, Franco. "The Monetarist Controversy, or Should We Forsake Stabilization Policies?" \textit{American Economic Review} (Março 1977): 1-19.
    \item[{[22]}] Mortensen, Dale T. "A Theory of Wage and Employment Dynamics," em [25].
    \item[{[23]}] Muench, T.; Rolnick, A.; Wallace, N.; e Weiler, W. "Tests for Structural Change and Prediction Intervals for the Reduced Forms of Two Structural Models of the U.S.: The FRB-MIT and Michigan Quarterly Models." \textit{Annals of Economic and Social Measurement} 313 (1974).
    \item[{[24]}] Okun, Arthur, e Perry, George L., eds. \textit{Brookings Papers on Economic Activity}, 1973, Vol. 3. Observações atribuídas a Lawrence Klein, p. 644.
    \item[{[25]}] Phelps, E. S. et al. \textit{Microeconomic Foundations of Employment and Inflation Theory}. New York: Norton, 1970.
    \item[{[26]}] \underline{\hspace{1cm}} e Taylor, John. "Stabilizing Powers of Monetary Policy under Rational Expectations." \textit{Journal of Political Economy} 85 (1977): 163-190.
    \item[{[27]}] Sargent, T. J. "A Classical Macroeconometric Model for the United States." \textit{Journal of Political Economy} (1976).
    \item[{[28]}] \underline{\hspace{1cm}}. "The Observational Equivalence of Natural and Unnatural Rate Theories of Macroeconomics." \textit{Journal of Political Economy} (Junho 1976).
    \item[{[29]}] \underline{\hspace{1cm}}. "Estimation of Dynamic Labor Demand Schedules Under Rational Expectations." \textit{Journal of Political Economy} (Dezembro 1978).
    \item[{[30]}] \underline{\hspace{1cm}} e Wallace, Neil. "Rational Expectations, the Optimal Monetary Instrument, and the Optimal Money Supply Rule." \textit{Journal of Political Economy} 83 (1975): 241-54.
    \item[{[31]}] \underline{\hspace{1cm}} e Sims, C.A. "Business Cycle Modeling Without Pretending to Have Too Much A Priori Economic Theory." Em \textit{New Methods in Business Cycle Research}, editado por C. Sims. Federal Reserve Bank of Minneapolis, 1977.
    \item[{[32]}] Simons, Henry C. "Rules Versus Authorities in Monetary Policy." \textit{Journal of Political Economy} 44 (1936): 1-30.
    \item[{[33]}] Sims, C. A. "Macroeconomics and Reality." \textit{Econometrica} (no prelo).
    \item[{[34]}] \underline{\hspace{1cm}}. "Money, Income, and Causality." \textit{American Economic Review} (Setembro 1972).
    \item[{[35]}] \underline{\hspace{1cm}}, ed. \textit{New Methods in Business Cycle Research: Proceedings from a Conference}. Federal Reserve Bank of Minneapolis, Outubro, 1977.
    \item[{[36]}] Sonnenschein, Hugo. "Do Walras' Identity and Continuity Characterize the Class of Community Excess Demand Functions?" \textit{Journal of Economic Theory} 6 (1973): 345-354.
    \item[{[37]}] Tobin, James. "Money Wage Rates and Employment." Em \textit{The New Economics}, editado por S. Harris. Clifton, New Jersey: Augustus Kelley, 1965.
    \item[{[38]}] \underline{\hspace{1cm}}. "How Dead is Keynes?" \textit{Economic Inquiry}, (Outubro 1977) 15: 459-68.
\end{enumerate}

\end{document}